
Automated decision-making systems,
based on ML models, 
are now trusted to make decisions of legal consequence, such as extending lines of credit, 
matching employees with employers, etc.
These systems are typically built atop supervised ML models.
In practice, the ML models 
do not take any actions themselves:
they simply estimate the conditional probability 
of a label given some features.
As a concrete example,
a label might be a binary indicator $\left \lbrace 0, 1 \right\rbrace$ 
of whether a loan defaults, 
and the features might be attributes 
related to a loan applicant's financial history.
Decisions are typically made by feeding the conditional probabilities into some decision rule. 

In the simplest systems, the decision rule consists of a threshold applied to the prediction.
For example, if $\hat P(\textit{default}|\mathbf{x}) > .2 \Rightarrow \textit{reject loan}$.
In many cases, while model outputs are generated programatically,
decisions are made by humans. 
For example recidivism scores 
-- predicted probabilities that an individual
will be rearrested after being released --
are taken into account as one among many criteria 
by judges when making decisions about bail, parole and sentencing.

A major concern is that a decision-making system based on an ML model
might exhibit behavior 
that is unlawful with respect to protected characteristics.
If the model explicitly incorporates a protected characteristic as input,
this amounts to disparate treatment \cite{CivilRights1964, barocas2016big, zafar2017fairness}.
An automated ML-based system can also exhibit disparate impact,
which could arise via several possible mechanisms: 
(i) The choice of the target may be arbitrarily chosen 
among several that are all loose surrogates 
for the real quantity of interest (e.g.~credit-worthiness).
Any one particular choice might then benefit or disadvantage a given group;
(ii) The targets themselves might reflect patterns of historical prejudice.
For example, if members of the dominant group 
are more likely to keep their jobs during a layoff, 
then they may be more likely to repay their loans;
(iii) If a group is grossly under-represented in the dataset, 
that could potentially lead to a model 
that is not as accurate when evaluating members of that group.
In some of the literature on fairness in algorithmic decision-making, 
use of the protected characteristic is called \emph{direct discrimination},
while points (i-iii) are sometime described as \emph{indirect discrimination} \cite{pedreshi2008discrimination,kamishima2011fairness}.


Note that the behavior of an ML model 
depends on the dataset used for training, 
and thus these mechanisms could induce disparate impact absent a data scientist's knowledge. 
As Barocas and Selbst note, 
``honest  attempts  to  certify the absence of prejudice on the
part of those involved in the data
mining process may wrongly confer the 
imprimatur of impartiality on the resulting decisions'' \cite{barocas2016big}.
These issues have come to light starting with the ground-breaking paper \cite{friedman1996bias},
which recognized the ability of computer systems to disparately affect protected groups.
Work specifically focused on discrimination owing to data-mining
became more common following \cite{pedreshi2008discrimination}. 
Recently, several prominent cases of disparate impact 
garnered widespread media attention.
For example, a report on ProPublica 
suggested that automated recidivism risk scores, 
used in sentencing and bail decisions, 
show racial bias \cite{MachineBias20}.
Following the organization of workshops 
and conferences addressing related topics,
research into algorithmic methods for mitigating
disparate impact has accelerated.
